\documentclass[11pt,a4paper]{article}
\usepackage{fontspec,amsmath,amssymb,amsthm,microtype}
\setmainfont{lmroman10-regular.otf}[BoldFont=lmroman10-bold.otf,ItalicFont=lmroman10-italic.otf,BoldItalicFont=lmroman10-bolditalic.otf]
\usepackage{xcolor}
\usepackage[most]{tcolorbox}
\usepackage{mathrsfs}
\providecommand{\lt}{<}
\providecommand{\gt}{>}
\setmonofont{DejaVuSansMono.ttf}[BoldFont=DejaVuSansMono-Bold.ttf,ItalicFont=DejaVuSansMono-Oblique.ttf,BoldItalicFont=DejaVuSansMono-BoldOblique.ttf,Scale=0.86]
\newfontfamily\leanfallback{STIXTwoMath-Regular.otf}[Scale=0.86]
\definecolor{leanrule}{gray}{0.62}
\definecolor{leanback}{gray}{0.962}
\definecolor{leantitle}{gray}{0.22}
\definecolor{leanmuted}{gray}{0.38}
\newtcolorbox{leanbox}[1]{enhanced,breakable,colback=leanback,colframe=leanrule,
  boxrule=0.4pt,arc=0.6mm,left=2.4mm,right=2.4mm,top=1.8mm,bottom=1.8mm,
  before skip=0.8em,after skip=1em,fonttitle=\small\bfseries,coltitle=leantitle,
  colbacktitle=leanback,titlerule=0.3pt,title={#1}}
\newlength{\leancharwidth}
\newenvironment{leancode}{\par\smallskip\begingroup\ttfamily\footnotesize\raggedright
  \settowidth{\leancharwidth}{M}\setlength{\parskip}{0pt}\setlength{\parindent}{0pt}
  \frenchspacing}{\par\endgroup\smallskip}
\newcommand{\leanline}[2]{\hangindent=\dimexpr#1\leancharwidth+4\leancharwidth\relax
  \hangafter=1\noindent\hspace*{#1\leancharwidth}#2\par}
\newcommand{\leansource}[1]{\par{\footnotesize\color{leanmuted}#1}\par}
\newcommand{\leanlabel}[1]{\par\smallskip\noindent{\small\bfseries #1}\par}
\usepackage[margin=25mm]{geometry}
\usepackage[unicode,hidelinks]{hyperref}
\hypersetup{pdfauthor={LangevinLean Project}}
\setlength{\emergencystretch}{3em}
\setlength{\parskip}{0.4em}
\setlength{\parindent}{0pt}
\allowdisplaybreaks[2]
\raggedbottom
\title{\texorpdfstring{The genuine Hessian and its segment average}{The genuine Hessian and its segment average}}
\author{LangevinLean Project}
\date{}
\begin{document}\maketitle
Draft mathematical exposition; these notes may contain mistakes.


The curvature operator is derived from the actual potential. For \(V\in C^2(H)\), set \(A(x)=D(\nabla V)(x)\) and \(\overline A(x,y)=\int_0^1 A(y+t(x-y))\,dt\). The construction proves symmetry, sharp curvature bounds, continuity and the exact gradient secant identity.


\section{\texorpdfstring{The actual second derivative}{The actual second derivative}}


\textbf{Claim.} The gradient is continuously differentiable, \(\langle A(x)u,v\rangle=D^2V(x)[u,v]\), and \(A(x)\) is symmetric and continuous in \(x\).



\textbf{Proof.} Identify the gradient with the inverse Riesz image of the first derivative. Differentiation through this continuous linear identification gives the displayed bilinear form. Equality of mixed second derivatives gives symmetry. Continuous second differentiability gives continuity of the derivative of the gradient.

\begin{leanbox}{Lean: \texttt{contDiff\_\allowbreak{}gradient\_\allowbreak{}of\_\allowbreak{}two}, \texttt{hessianOperator}, \texttt{inner\_\allowbreak{}hessianOperator} and 2 more}
\leansource{\texttt{contDiff\_\allowbreak{}gradient\_\allowbreak{}of\_\allowbreak{}two} \quad WassersteinMixing/\allowbreak{}Curvature.lean, lines 24–27}
\begin{leancode}
\leanline{0}{theorem\ contDiff\_gradient\_of\_two\ (V\ :\ H\ →\ ℝ)\ (hV\ :\ ContDiff\ ℝ\ 2\ V)\ :}
\leanline{4}{ContDiff\ ℝ\ 1\ (gradient\ V)}
\end{leancode}
\leansource{\texttt{hessianOperator} \quad WassersteinMixing/\allowbreak{}Curvature.lean, lines 30–30}
\begin{leancode}
\leanline{0}{def\ hessianOperator\ (V\ :\ H\ →\ ℝ)\ (x\ :\ H)\ :\ H\ →L[ℝ]\ H\ :=\ fderiv\ ℝ\ (gradient\ V)\ x}
\end{leancode}
\leansource{\texttt{inner\_\allowbreak{}hessianOperator} \quad WassersteinMixing/\allowbreak{}Curvature.lean, lines 33–43}
\begin{leancode}
\leanline{0}{theorem\ inner\_hessianOperator\ (V\ :\ H\ →\ ℝ)\ (hV\ :\ ContDiff\ ℝ\ 2\ V)\ (x\ u\ v\ :\ H)\ :}
\leanline{4}{inner\ ℝ\ (hessianOperator\ V\ x\ u)\ v\ =\ fderiv\ ℝ\ (fderiv\ ℝ\ V)\ x\ u\ v}
\end{leancode}
\leansource{\texttt{hessianOperator\_\allowbreak{}symmetric} \quad WassersteinMixing/\allowbreak{}Curvature.lean, lines 46–49}
\begin{leancode}
\leanline{0}{theorem\ hessianOperator\_symmetric\ (V\ :\ H\ →\ ℝ)\ (hV\ :\ ContDiff\ ℝ\ 2\ V)\ (x\ u\ v\ :\ H)\ :}
\leanline{4}{inner\ ℝ\ (hessianOperator\ V\ x\ u)\ v\ =\ inner\ ℝ\ u\ (hessianOperator\ V\ x\ v)}
\end{leancode}
\leansource{\texttt{hessianOperator\_\allowbreak{}continuous} \quad WassersteinMixing/\allowbreak{}Curvature.lean, lines 52–54}
\begin{leancode}
\leanline{0}{theorem\ hessianOperator\_continuous\ (V\ :\ H\ →\ ℝ)\ (hV\ :\ ContDiff\ ℝ\ 2\ V)\ :}
\leanline{4}{Continuous\ (hessianOperator\ V)}
\end{leancode}
\leanlabel{How the Lean proof follows the argument}
\texttt{cont\allowbreak{}Diff\_\allowbreak{}gradient\_\allowbreak{}of\_\allowbreak{}two} differentiates the Riesz representation. \texttt{hessian\allowbreak{}Operator} is the actual Fréchet derivative of the gradient. \texttt{inner\_\allowbreak{}hessian\allowbreak{}Operator} identifies its form with the second derivative; symmetry follows from the second-derivative symmetry theorem.
\end{leanbox}

\section{\texorpdfstring{Pointwise curvature bounds}{Pointwise curvature bounds}}


\textbf{Claim.} Under the potential contract, \(\|A(x)\|\le\beta\) and

\[\alpha\|v\|^2\le\langle v,A(x)v\rangle\le\beta\|v\|^2.\]

\textbf{Proof.} A derivative of a \(\beta\)-Lipschitz map has norm at most \(\beta\), giving the upper estimate by Cauchy–Schwarz. Strong convexity makes the gradient \(\alpha\)-strongly monotone. Apply that inequality to \(x+tv\) and \(x\), divide by the positive displacement parameter, and pass to the directional derivative as \(t\downarrow0\) to obtain the lower bound.

\begin{leanbox}{Lean: \texttt{hessianOperator\_\allowbreak{}norm\_\allowbreak{}le}, \texttt{hessianOperator\_\allowbreak{}lower}, \texttt{hessianOperator\_\allowbreak{}upper}}
\leansource{\texttt{hessianOperator\_\allowbreak{}norm\_\allowbreak{}le} \quad WassersteinMixing/\allowbreak{}Curvature.lean, lines 57–59}
\begin{leancode}
\leanline{0}{theorem\ hessianOperator\_norm\_le\ (P\ :\ Parameters)\ (V\ :\ H\ →\ ℝ)}
\leanline{4}{(hV\ :\ PotentialAssumptions\ P\ V)\ (x\ :\ H)\ :\ ‖hessianOperator\ V\ x‖\ ≤\ P.beta}
\end{leancode}
\leansource{\texttt{hessianOperator\_\allowbreak{}lower} \quad WassersteinMixing/\allowbreak{}Curvature.lean, lines 62–85}
\begin{leancode}
\leanline{0}{theorem\ hessianOperator\_lower\ (P\ :\ Parameters)\ (V\ :\ H\ →\ ℝ)}
\leanline{4}{(hV\ :\ PotentialAssumptions\ P\ V)\ (x\ v\ :\ H)\ :}
\leanline{4}{P.alpha\ *\ ‖v‖\ \textasciicircum{}\ 2\ ≤\ inner\ ℝ\ v\ (hessianOperator\ V\ x\ v)}
\end{leancode}
\leansource{\texttt{hessianOperator\_\allowbreak{}upper} \quad WassersteinMixing/\allowbreak{}Curvature.lean, lines 88–94}
\begin{leancode}
\leanline{0}{theorem\ hessianOperator\_upper\ (P\ :\ Parameters)\ (V\ :\ H\ →\ ℝ)}
\leanline{4}{(hV\ :\ PotentialAssumptions\ P\ V)\ (x\ v\ :\ H)\ :}
\leanline{4}{inner\ ℝ\ v\ (hessianOperator\ V\ x\ v)\ ≤\ P.beta\ *\ ‖v‖\ \textasciicircum{}\ 2}
\end{leancode}
\leanlabel{How the Lean proof follows the argument}
\texttt{hessian\allowbreak{}Operator\_\allowbreak{}norm\_\allowbreak{}le} uses the derivative bound of a Lipschitz map. \texttt{hessian\allowbreak{}Operator\_\allowbreak{}lower} takes a one-sided slope limit of the strong-gradient-monotonicity inequality. \texttt{hessian\allowbreak{}Operator\_\allowbreak{}upper} combines the norm estimate with the real inner-product bound.
\end{leanbox}


Earlier results used here: \cite[Section~3.2.3]{LangevinLean2026ULASC}.

\section{\texorpdfstring{The exact averaged secant operator}{The exact averaged secant operator}}


\textbf{Claim.} The actual operator integral exists and satisfies

\[\overline A(x,y)(x-y)=\nabla V(x)-\nabla V(y),\qquad \alpha I\preceq\overline A(x,y)\preceq\beta I.\]

\textbf{Proof.} Its integrand is continuous on a compact interval. Evaluation and inner-product functionals commute with the Bochner integral. Integrating the pointwise symmetry and form bounds therefore preserves them, and the norm bound is preserved because the interval has length one. The derivative of the gradient along the segment is \(A(y+t(x-y))(x-y)\); the fundamental theorem gives the secant identity. Joint continuity in the endpoints follows by dominated convergence, with the uniform bound \(\beta\).

\begin{leanbox}{Lean: \texttt{averagedHessian}, \texttt{averagedHessian\_\allowbreak{}curve\_\allowbreak{}continuous}, \texttt{averagedHessian\_\allowbreak{}apply} and 7 more}
\leansource{\texttt{averagedHessian} \quad WassersteinMixing/\allowbreak{}Curvature.lean, lines 97–98}
\begin{leancode}
\leanline{0}{def\ averagedHessian\ (V\ :\ H\ →\ ℝ)\ (x\ y\ :\ H)\ :\ H\ →L[ℝ]\ H\ :=}
\leanline{2}{∫\ t\ in\ (0\ :\ ℝ)..1,\ hessianOperator\ V\ (y\ +\ t\ •\ (x\ -\ y))}
\end{leancode}
\leansource{\texttt{averagedHessian\_\allowbreak{}curve\_\allowbreak{}continuous} \quad WassersteinMixing/\allowbreak{}Curvature.lean, lines 100–102}
\begin{leancode}
\leanline{0}{theorem\ averagedHessian\_curve\_continuous\ (V\ :\ H\ →\ ℝ)\ (hV\ :\ ContDiff\ ℝ\ 2\ V)\ (x\ y\ :\ H)\ :}
\leanline{4}{Continuous\ (fun\ t\ :\ ℝ\ ↦\ hessianOperator\ V\ (y\ +\ t\ •\ (x\ -\ y)))}
\end{leancode}
\leansource{\texttt{averagedHessian\_\allowbreak{}apply} \quad WassersteinMixing/\allowbreak{}Curvature.lean, lines 105–108}
\begin{leancode}
\leanline{0}{theorem\ averagedHessian\_apply\ (V\ :\ H\ →\ ℝ)\ (hV\ :\ ContDiff\ ℝ\ 2\ V)\ (x\ y\ v\ :\ H)\ :}
\leanline{4}{averagedHessian\ V\ x\ y\ v\ =\ ∫\ t\ in\ (0\ :\ ℝ)..1,\ hessianOperator\ V\ (y\ +\ t\ •\ (x\ -\ y))\ v}
\end{leancode}
\leansource{\texttt{averagedHessian\_\allowbreak{}inner} \quad WassersteinMixing/\allowbreak{}Curvature.lean, lines 111–116}
\begin{leancode}
\leanline{0}{theorem\ averagedHessian\_inner\ (V\ :\ H\ →\ ℝ)\ (hV\ :\ ContDiff\ ℝ\ 2\ V)\ (x\ y\ u\ v\ :\ H)\ :}
\leanline{4}{inner\ ℝ\ u\ (averagedHessian\ V\ x\ y\ v)\ =}
\leanline{6}{∫\ t\ in\ (0\ :\ ℝ)..1,\ inner\ ℝ\ u\ (hessianOperator\ V\ (y\ +\ t\ •\ (x\ -\ y))\ v)}
\end{leancode}
\leansource{\texttt{averagedHessian\_\allowbreak{}symmetric} \quad WassersteinMixing/\allowbreak{}Curvature.lean, lines 119–125}
\begin{leancode}
\leanline{0}{theorem\ averagedHessian\_symmetric\ (V\ :\ H\ →\ ℝ)\ (hV\ :\ ContDiff\ ℝ\ 2\ V)\ (x\ y\ u\ v\ :\ H)\ :}
\leanline{4}{inner\ ℝ\ (averagedHessian\ V\ x\ y\ u)\ v\ =\ inner\ ℝ\ u\ (averagedHessian\ V\ x\ y\ v)}
\end{leancode}
\leansource{\texttt{averagedHessian\_\allowbreak{}norm\_\allowbreak{}le} \quad WassersteinMixing/\allowbreak{}Curvature.lean, lines 128–133}
\begin{leancode}
\leanline{0}{theorem\ averagedHessian\_norm\_le\ (P\ :\ Parameters)\ (V\ :\ H\ →\ ℝ)}
\leanline{4}{(hV\ :\ PotentialAssumptions\ P\ V)\ (x\ y\ :\ H)\ :\ ‖averagedHessian\ V\ x\ y‖\ ≤\ P.beta}
\end{leancode}
\leansource{\texttt{averagedHessian\_\allowbreak{}lower} \quad WassersteinMixing/\allowbreak{}Curvature.lean, lines 136–147}
\begin{leancode}
\leanline{0}{theorem\ averagedHessian\_lower\ (P\ :\ Parameters)\ (V\ :\ H\ →\ ℝ)}
\leanline{4}{(hV\ :\ PotentialAssumptions\ P\ V)\ (x\ y\ v\ :\ H)\ :}
\leanline{4}{P.alpha\ *\ ‖v‖\ \textasciicircum{}\ 2\ ≤\ inner\ ℝ\ v\ (averagedHessian\ V\ x\ y\ v)}
\end{leancode}
\leansource{\texttt{averagedHessian\_\allowbreak{}upper} \quad WassersteinMixing/\allowbreak{}Curvature.lean, lines 150–156}
\begin{leancode}
\leanline{0}{theorem\ averagedHessian\_upper\ (P\ :\ Parameters)\ (V\ :\ H\ →\ ℝ)}
\leanline{4}{(hV\ :\ PotentialAssumptions\ P\ V)\ (x\ y\ v\ :\ H)\ :}
\leanline{4}{inner\ ℝ\ v\ (averagedHessian\ V\ x\ y\ v)\ ≤\ P.beta\ *\ ‖v‖\ \textasciicircum{}\ 2}
\end{leancode}
\leansource{\texttt{averagedHessian\_\allowbreak{}secant} \quad WassersteinMixing/\allowbreak{}Curvature.lean, lines 160–170}
\begin{leancode}
\leanline{0}{theorem\ averagedHessian\_secant\ (V\ :\ H\ →\ ℝ)\ (hV\ :\ ContDiff\ ℝ\ 2\ V)\ (x\ y\ :\ H)\ :}
\leanline{4}{averagedHessian\ V\ x\ y\ (x\ -\ y)\ =\ gradient\ V\ x\ -\ gradient\ V\ y}
\end{leancode}
\leansource{\texttt{averagedHessian\_\allowbreak{}continuous} \quad WassersteinMixing/\allowbreak{}Curvature.lean, lines 174–189}
\begin{leancode}
\leanline{0}{theorem\ averagedHessian\_continuous\ (P\ :\ Parameters)\ (V\ :\ H\ →\ ℝ)}
\leanline{4}{(hV\ :\ PotentialAssumptions\ P\ V)\ :}
\leanline{4}{Continuous\ (fun\ p\ :\ H\ ×\ H\ ↦\ averagedHessian\ V\ p.1\ p.2)}
\end{leancode}
\leanlabel{How the Lean proof follows the argument}
\texttt{averaged\allowbreak{}Hessian} is an operator-valued interval integral. The application and pairing lemmas justify moving bounded linear functionals inside it. The secant theorem differentiates the actual gradient curve. Joint continuity uses the proved uniform Hessian bound as a dominated-convergence envelope.
\end{leanbox}

\begin{thebibliography}{ULA-LSI}
\bibitem[ULA-SC]{LangevinLean2026ULASC}
LangevinLean Project. \emph{Unadjusted Langevin under Strong Convexity}. LL-01, 2026. Draft report, content revision d34ead06987c. \url{https://odakyildiz.github.io/LangevinLean/docs/pdf/UnadjustedLangevin/FullProof.pdf}.
\end{thebibliography}

\end{document}
